% paper3_mu0_dimensional.tex
% Supplemental Material: Dimensional uniqueness of mu_0 = H_0/c
% Supplement to Paper 3 (Dark Matter as Khronon Condensate)
% Author: Sheng-Kai Huang
% Compile: pdflatex paper3_mu0_dimensional.tex && pdflatex paper3_mu0_dimensional.tex

\documentclass[prd,twocolumn,superscriptaddress,longbibliography,floatfix]{revtex4-2}

\usepackage{amsmath}
\usepackage{amssymb}
\usepackage{amsfonts}
\usepackage{amsthm}
\usepackage{bm}
\usepackage{hyperref}
\usepackage{booktabs}
\usepackage{xcolor}

\newtheorem{theorem}{Theorem}
\newtheorem{proposition}[theorem]{Proposition}
\newtheorem{corollary}[theorem]{Corollary}
\newtheorem*{remark}{Remark}

\newcommand{\Hzero}{H_{0}}
\newcommand{\muzero}{\mu_{0}}
\newcommand{\rhocrit}{\rho_{\mathrm{crit}}}
\newcommand{\OmegaDM}{\Omega_{\mathrm{DM}}}
\newcommand{\MPl}{M_{\mathrm{Pl}}}
\newcommand{\known}{\textsc{[known]}}
\newcommand{\derived}{\textsc{[derived]}}
\newcommand{\assumed}{\textsc{[assumed]}}
\newcommand{\new}{\textsc{[new synthesis]}}

\begin{document}

\title{Supplemental Material: Dimensional Uniqueness of\\
$\muzero = \Hzero/c$ in Khronon Ghost Condensation}

\author{Sheng-Kai Huang}
\email{akai@fawstudio.com}
\affiliation{Independent Researcher}

\date{\today}

\begin{abstract}
We show that the cosmological scale appearing in the
Blanchet--Skordis kinetic function $K(Q) = \mu^{2}(Q-1)^{2}$
is uniquely fixed, up to an $O(1)$ prefactor, to be
$\muzero = \Hzero/c$.  Two independent physical
constraints---(i) the requirement that the condensate
energy density reproduce $\rhocrit \sim \Hzero^{2}/G$
without a hierarchy, and (ii) the requirement that
$\muzero$ is the IR scale where Petz retrodictability
fails---both single out $\Hzero/c$ as the unique
dimensionally consistent choice.  Other cosmological
scales ($\sqrt{\Lambda}$, $G\rho_{\mathrm{m}}/c^{4}$,
particle masses) either reduce to $\Hzero/c$ at $z = 0$
or fail one of the two constraints.  The numerical
prefactor $1$ versus $1/(2\pi)$ versus $2$ is not fixed
by dimensional analysis; in the main text the
combination $a_{0} = c^{2}\muzero/(2\pi)$ that appears
in the galactic sector singles out the $1/(2\pi)$
factor through the de~Broglie--like spatial Fourier
normalization.
\end{abstract}

\maketitle

%=======================================================================
\section{Setup}
\label{sec:setup}
%=======================================================================

\known\ The Blanchet--Skordis (BS) Khronon
action~\cite{BS2024,BS2025} contains a kinetic function
of the timelike scalar
\begin{equation}
K(Q) = \mu^{2}(Q-1)^{2} + \cdots,
\label{eq:KQ}
\end{equation}
where $Q = -g^{\mu\nu}\nabla_{\mu}\tau\nabla_{\nu}\tau$
is dimensionless and $\mu$ has dimension $[\mathrm{length}]^{-1}$
(in units with $\hbar = 1$, equivalently
$[\mathrm{mass}]$ in natural units).  Ghost condensation
sits at the minimum $Q = 1$.  The cosmological background
of Paper~3 is at $Q_{0} = 1 + \delta_{0}$ with
$\delta_{0} \approx 0.34$.

The dark matter density follows from
\begin{equation}
\rho_{K} = c^{2}\,\frac{\mu^{2}}{8\pi G}\,\delta_{0}(2+\delta_{0}),
\label{eq:rhoK}
\end{equation}
so that the dimensionless density parameter is
\begin{equation}
\OmegaDM = \frac{\rho_{K}}{\rhocrit}
= \frac{c^{2}\mu^{2}}{\Hzero^{2}}\,
\frac{\delta_{0}(2+\delta_{0})}{3},
\label{eq:OmegaDM_general}
\end{equation}
using $\rhocrit = 3\Hzero^{2}/(8\pi G)$.

The question is: \emph{what fixes the scale $\mu$?}
The BS papers leave $\mu$ as a free parameter set by
matching to observations.  We argue here that one
\emph{specific} choice is selected by dimensional analysis
combined with the ghost-condensation IR structure---and
that this choice happens to reproduce $\OmegaDM$ within
$0.5\%$ once $\delta_{0}$ is fitted, with no further
hierarchy.

%=======================================================================
\section{Available scales}
\label{sec:scales}
%=======================================================================

\derived\ At cosmological scales today, the available
fundamental constants and parameters with dimension are
\begin{itemize}
\item the speed of light $c$
  ($[c] = \mathrm{m/s}$);
\item Newton's constant $G$
  ($[G] = \mathrm{m}^{3}\mathrm{kg}^{-1}\mathrm{s}^{-2}$);
\item the Hubble rate $\Hzero$
  ($[\Hzero] = \mathrm{s}^{-1}$);
\item particle masses $m_{i}$
  ($[m_{i}] = \mathrm{kg}$, with the lightest neutrino
  $m_{\nu} \gtrsim 10^{-2}\,\mathrm{eV}/c^{2}$);
\item the cosmological constant
  $\Lambda \sim \Hzero^{2}/c^{2}$
  (not independent of $\Hzero$);
\item the Planck mass
  $\MPl = (\hbar c/G)^{1/2}$ (with $\hbar$ included only
  for quantum gravity).
\end{itemize}

We restrict $\hbar$ to enter only via $\MPl$ (the
classical limit).  Then any inverse length must be
expressible as
\begin{equation}
\mu \;\propto\; \Hzero^{n}\,c^{m}\,G^{k}\,m_{i}^{\ell},
\label{eq:mu_ansatz}
\end{equation}
for real exponents $n,m,k,\ell$ satisfying
\begin{equation}
[\mathrm{m}^{-1}]
= [\mathrm{s}^{-n}]\cdot[\mathrm{m/s}]^{m}\cdot
  [\mathrm{m}^{3}/(\mathrm{kg}\,\mathrm{s}^{2})]^{k}\cdot
  [\mathrm{kg}]^{\ell},
\label{eq:dimcount}
\end{equation}
giving three constraints:
$3k + m = -1$ (length),
$-n - m - 2k = 0$ (time),
$-k + \ell = 0$ (mass).
This is a one-parameter family
$\{n,m,k,\ell\} = \{n,\,n-1,\,0,\,0\} + (\text{$G$-axis})$.

\textbf{Removing irrelevant axes.}  Ghost condensation
is an IR mechanism: the kinetic function $K(Q)$ describes
classical, long-wavelength dynamics of the Khronon, so
$\hbar$ should not appear, removing the $\MPl$ axis.
At cosmological scales the only relevant matter
density is $\rho_{\mathrm{m}}$; combinations such as
$G\rho_{\mathrm{m}}/c^{2}$ have dimension of $\Hzero^{2}/c^{2}$
on the FRW background by Friedmann's equation, so any
$G$-dependence collapses onto $\Hzero^{2}$.  Particle
masses $m_{i}$ are at energies $\gtrsim 10^{15}\Hzero$
above the cosmological IR scale and would introduce a
hierarchy of $\sim 10^{60}$ in $\OmegaDM$ if they entered
$\mu$ directly.

After these reductions, the only dimensionally consistent
scale linear in $\Hzero$ is
\begin{equation}
\boxed{\;\muzero \,=\, \xi\,\frac{\Hzero}{c}\,,\;}
\label{eq:mu0}
\end{equation}
with $\xi$ a dimensionless $O(1)$ constant.

%=======================================================================
\section{Fixing the prefactor}
\label{sec:prefactor}
%=======================================================================

\subsection{Constraint 1: matching $\rhocrit$}

\derived\ Substituting $\muzero = \xi\Hzero/c$
into~\eqref{eq:rhoK},
\begin{equation}
\rho_{K} = \frac{c^{2}}{8\pi G}\,\frac{\xi^{2}\Hzero^{2}}{c^{2}}\,
\delta_{0}(2+\delta_{0})
= \frac{\xi^{2}\Hzero^{2}}{8\pi G}\,\delta_{0}(2+\delta_{0})\,,
\label{eq:rhoK_xi}
\end{equation}
so that
\begin{equation}
\OmegaDM = \frac{\xi^{2}\,\delta_{0}(2+\delta_{0})}{3}\,.
\label{eq:OmegaDM_xi}
\end{equation}
With $\xi = 1$ and $\delta_{0} = 0.340$,
$\OmegaDM = 0.2652$, agreeing with Planck~2018
$\OmegaDM = 0.2665$~\cite{Planck2018} to $0.49\%$.
Any $\xi \neq 1$ requires a compensating shift in
$\delta_{0}$ (or vice versa); $\xi = 1$ is the
\emph{minimal} choice with no fine-tuning of the
prefactor against $\delta_{0}$.

\subsection{Constraint 2: IR retrodictability scale}

\new\ The Petz recovery framework of Paper~1
\cite{Huang2025paper1} identifies $\Hzero^{-1}$ as the
\emph{retrodictability horizon}: signals beyond this
scale leak from the observable Hilbert space, and the
recovery fidelity $F = e^{-\Sigma/2}$ degrades.
For Khronon condensation to encode the cosmological arrow
of time, the IR scale of the kinetic function $K(Q)$
must coincide with this horizon:
\begin{equation}
\mu^{-1} \,\sim\, c\Hzero^{-1} \;=\; d_{H},
\label{eq:horizon}
\end{equation}
the Hubble length.  This is satisfied by
$\muzero = \Hzero/c$ and only by that choice.
Multiplying or dividing by $2\pi$ shifts $\mu^{-1}$ by
$2\pi$, which corresponds to going from the wavelength
to the wavenumber---this is the source of the $1/(2\pi)$
in the galactic sector $a_{0} = c^{2}\muzero/(2\pi)$.

%=======================================================================
\section{Numerical evaluation}
\label{sec:numerical}
%=======================================================================

\derived\ With $\Hzero = 67.4\;\mathrm{km/s/Mpc}
= 2.184 \times 10^{-18}\;\mathrm{s}^{-1}$
\cite{Planck2018}:
\begin{align}
\muzero
&= \frac{\Hzero}{c}
= \frac{2.184 \times 10^{-18}\;\mathrm{s}^{-1}}
       {2.998 \times 10^{8}\;\mathrm{m/s}}\nonumber\\
&= 7.286 \times 10^{-27}\;\mathrm{m}^{-1}.
\label{eq:mu0_num}
\end{align}
The corresponding inverse length is
$\muzero^{-1} = 1.372 \times 10^{26}\;\mathrm{m}
\approx d_{H}$, the Hubble radius.

\textbf{Cross-check via $\rhocrit$.}
\begin{align}
\frac{c^{2}\muzero^{2}}{8\pi G}
&= \frac{c^{2}(\Hzero/c)^{2}}{8\pi G}
= \frac{\Hzero^{2}}{8\pi G}\nonumber\\
&= \frac{\rhocrit}{3}
= 2.844 \times 10^{-27}\;\mathrm{kg/m}^{3},
\label{eq:rho_check}
\end{align}
exactly $\rhocrit/3$.  The factor of $3$ is absorbed
into $\delta_{0}(2+\delta_{0})$, recovering
$\OmegaDM = 0.265$.

\textbf{Cross-check via $a_{0}$.}  The galactic sector
predicts the MOND scale
\begin{align}
a_{0} = \frac{c^{2}\muzero}{2\pi}
&= \frac{c\Hzero}{2\pi}
= \frac{(2.998 \times 10^{8})(2.184 \times 10^{-18})}{2\pi}\nonumber\\
&= 1.042 \times 10^{-10}\;\mathrm{m/s}^{2},
\label{eq:a0_check}
\end{align}
within $13\%$ of the empirical
$a_{0} \approx 1.2 \times 10^{-10}\;\mathrm{m/s}^{2}$
\cite{Milgrom1983}.  We emphasize this is suggestive, not
a derivation of the empirical $a_{0}$: the BS galactic
sector $J(Y)$ is a free function and $a_{0}$ enters as
its small-$Y$ asymptotic parameter.

%=======================================================================
\section{Why other scales fail}
\label{sec:fails}
%=======================================================================

We list candidate inverse-length scales and the constraint
they violate.

\begin{table}[b]
\caption{\label{tab:scales}Candidate cosmological scales
for $\mu$ and the constraint each fails.  Only
$\muzero = \Hzero/c$ passes both.}
\begin{ruledtabular}
\begin{tabular}{lll}
Scale & Numerical & Failure mode \\
\colrule
$\Hzero/c$ & $7.3 \times 10^{-27}\;\mathrm{m}^{-1}$
  & ---\;(passes both) \\
$2\Hzero/c$ & $1.5 \times 10^{-26}\;\mathrm{m}^{-1}$
  & $\OmegaDM = 1.06$ for fixed $\delta_{0}$ \\
$\Hzero/(2c)$ & $3.6 \times 10^{-27}\;\mathrm{m}^{-1}$
  & $\OmegaDM = 0.066$, $4\times$ low \\
$\sqrt{\Lambda}$ & $\sim \Hzero/c$
  & degenerate with $\Hzero/c$ \\
$\Hzero^{2}/(cM_{\nu})$ & $\sim 10^{-65}\;\mathrm{m}^{-1}$
  & $\OmegaDM \sim 10^{-77}$ \\
$1/\MPl$ & $1.6 \times 10^{35}\;\mathrm{m}^{-1}$
  & involves $\hbar$, IR/UV swap \\
$1/d_{\mathrm{LSS}}$ & $\sim \Hzero/c$
  & last-scattering, degenerate \\
\end{tabular}
\end{ruledtabular}
\end{table}

\textbf{$2\Hzero/c$ or $\Hzero/(2c)$:} simple
$O(1)$ rescalings change $\OmegaDM$ by a factor of
$\xi^{2}$, which would require shifting
$\delta_{0}$ to $0.158$ or $0.812$ respectively to
restore Planck consistency.  The preferred value
$\xi = 1$ keeps $\delta_{0}$ in its natural range
$\delta_{0} \sim O(0.1)$, with no tuning between
prefactor and amplitude.

\textbf{$\sqrt{\Lambda}$:} numerically
$\sqrt{\Lambda} \approx 1.27 \Hzero/c$ today, but
$\Lambda$ is constant and $\Hzero$ runs.  At
recombination, $\sqrt{\Lambda}$ is unchanged while
$H(z=1100)/c$ is $\sim 10^{5}\;\Hzero/c$; the two
scales diverge in the past.  $\sqrt{\Lambda}$ would
freeze the IR scale of $K(Q)$ at the present-day
horizon, giving the wrong scaling of $\rho_{K}(z)$.
The running coupling $\mu(a) = H(a)/c$ is
discussed in the companion supplement
\emph{paper3\_running\_coupling}.

\textbf{$\Hzero^{2}/(cM_{\nu})$ or
$1/\MPl$:} introduce a $10^{60}$--$10^{120}$ hierarchy
between $\mu$ and the observed $\OmegaDM$, requiring
fine-tuning at the cosmological constant level.

\textbf{Particle masses:} any $\mu \propto m_{i}c/\hbar$
sets the IR scale to a Compton wavelength, far below
the Hubble horizon, and is incompatible with the ghost
condensation IR structure.

%=======================================================================
\section{Connection to Petz recovery (Paper~1)}
\label{sec:petz}
%=======================================================================

\new\ Paper~1 establishes that $\Hzero$ is the natural
retrodictability scale of the universe: the entropy
production rate per Hubble time satisfies
$\dot{\Sigma}/\Hzero = O(1)$ on cosmological backgrounds.
The Khronon $Q$ is the order parameter of this
retrodictability, with $Q - 1$ measuring departure from
the time-symmetric vacuum.  In this picture
$\muzero = \Hzero/c$ is not an independent parameter
but the IR scale at which recovery fidelity
$F = e^{-\Sigma/2}$ is order unity over the observable
universe.

The factor of $2\pi$ between $\muzero$ and the galactic
$a_{0}$ scale arises because $a_{0}$ enters the
galactic sector through a spatial Fourier kernel
(de~Broglie--like normalization), while $\muzero$
enters the cosmological sector through the temporal
mass scale of $K(Q)$.  Both encode the same horizon,
viewed from the kinetic and gradient sides of the
Khronon action.

%=======================================================================
\section{Summary}
\label{sec:summary}
%=======================================================================

We have shown that $\muzero = \Hzero/c$ is the unique
inverse-length scale built from cosmological constants
that (i) does not invoke $\hbar$ or particle masses,
(ii) reproduces $\rhocrit$ to within the
$\delta_{0}(2+\delta_{0})/3$ factor without a
hierarchy, and (iii) coincides with the Petz
retrodictability horizon.  The
$O(1)$ prefactor $\xi = 1$ is selected by the absence
of a tuning between the kinetic scale and $\delta_{0}$;
the $1/(2\pi)$ factor in $a_{0}$ is the spatial
Fourier dual of the same scale.  No free parameter
remains in the cosmological sector beyond the boundary
condition $\delta_{0}$, and even that is fixed within
$2\sigma$ by Planck once $\xi = 1$ is imposed.

\begin{thebibliography}{99}

\bibitem{BS2024}
L.~Blanchet and C.~Skordis,
``Dark matter as a Khronon condensate,''
JCAP \textbf{11}, 040 (2024);
\href{https://arxiv.org/abs/2404.06584}{arXiv:2404.06584}.

\bibitem{BS2025}
L.~Blanchet and C.~Skordis,
``Khronon-Tensor theory reproducing MOND and the cosmological model,''
(2025);
\href{https://arxiv.org/abs/2507.00912}{arXiv:2507.00912}.

\bibitem{Planck2018}
Planck Collaboration,
``Planck 2018 results. VI. Cosmological parameters,''
Astron.\ Astrophys.\ \textbf{641}, A6 (2020);
\href{https://arxiv.org/abs/1807.06209}{arXiv:1807.06209}.

\bibitem{Milgrom1983}
M.~Milgrom,
``A modification of the Newtonian dynamics as a possible
alternative to the hidden mass hypothesis,''
Astrophys.\ J.\ \textbf{270}, 365 (1983).

\bibitem{Huang2025paper1}
S.-K.~Huang,
``The arrow of time as Petz recovery failure:
a unifying framework,''
(2025);
DOI: 10.5281/zenodo.14775757.

\bibitem{AHCLM2004}
N.~Arkani-Hamed, H.-C.~Cheng, M.~A.~Luty,
and S.~Mukohyama,
``Ghost condensation and a consistent infrared
modification of gravity,''
JHEP \textbf{05}, 074 (2004);
\href{https://arxiv.org/abs/hep-th/0312099}{arXiv:hep-th/0312099}.

\end{thebibliography}

\end{document}
